\section{Sustavi automatizirane vožnje i percepcija okoline}
\label{sec:percepcija}

Prije razmatranja normi koje osiguravaju sigurnost ADS-a korisno je razumjeti tehničku građu takvog sustava i, posebno, perceptivni lanac. Upravo u percepciji i donošenju odluka nastaje većina opasnosti kojima se bavi SOTIF, pa ovo poglavlje uspostavlja tehničku osnovu za poglavlja koja slijede.

\subsection{Razine automatizacije prema SAE J3016}

Norma SAE~J3016 \cite{saej3016} definira šest razina automatizacije vožnje. Na razinama 0--2 ljudski vozač i dalje nadzire dinamički zadatak vožnje, dok od razine~3 nadalje sustav samostalno izvodi vožnju unutar svoje operativne domene. Ova razlika ima izravne sigurnosne posljedice: na razinama 4 i 5 ne postoji ljudski vozač kao posljednja sigurnosna mreža, što iz temelja mijenja pretpostavke klasične funkcijske sigurnosti (vidi potpoglavlje~\ref{subsec:upravljivost}).

Primjeri sustava prisutnih na cesti ilustriraju ovu ljestvicu \cite{hrgetic2025uvod}: sustavi razine~2 uključuju Tesla Autopilot, Ford BlueCruise i GM SuperCruise; razinu~3 predstavlja Mercedes-Benz Drive Pilot; dok je primjer razine~4 Waymo robotaxi sustav. Vozilo pune autonomije (SAE~5) zamišlja se bez volana i pedala, potpuno prilagođeno udobnosti putnika.

\subsection{Arhitektura percepcijsko-upravljačkog lanca}

Zadaci percepcije autonomnog vozila obuhvaćaju percepciju okoline (senzorsko opažanje i interpretacija okruženja), procjenu vlastitog gibanja vozila, detekciju i klasifikaciju objekata, razumijevanje scene (struktura ceste, prometne trake, znakovi), praćenje dinamičkih objekata i predikciju njihovih trajektorija te lokalizaciju i mapiranje \cite{suligoj2025percepcija1}. Sva senzorska mjerenja obrađuju se i fuzioniraju kako bi se procijenila prostorna struktura okoline, pri čemu se prostor kontinuirano kategorizira kao zauzet (prepreke), slobodan (prohodne regije) ili nepoznat (područja bez pouzdanih mjerenja).

Na razini cjelovite arhitekture razlikuju se dva pristupa preslikavanju percepcije u akciju \cite{suligoj2025percepcija2}:
\begin{itemize}
  \item \textbf{Modularni pristup} (percepcija--planiranje--upravljanje): eksplicitni lanac u kojem se slika pretvara u strukturiranu reprezentaciju scene (trake, objekti, slobodan prostor), a zatim u odluku i akciju. Međurezultati su interpretabilni i provjerljivi, što olakšava otklanjanje pogrešaka i validaciju, no javlja se propagacija pogrešaka između modula;
  \item \textbf{End-to-end pristup} (slika--akcija): jedinstvena neuronska mreža zajednički uči percepciju i upravljanje bez eksplicitnih međureprezentacija. Arhitektura je jednostavnija i dobro skalira s podacima, ali je slabije interpretabilna te ju je teže validirati i sigurnosno potvrditi.
\end{itemize}
Niska interpretabilnost i otežana validacija end-to-end sustava izravno su povezane s izazovima koje SOTIF i norme za sigurnost umjetne inteligencije nastoje adresirati.

\subsection{Senzorske modalnosti}

Autonomna vozila koriste više senzorskih tehnologija koje se razlikuju po fizikalnom načelu i modalitetu podataka, ali sve funkcioniraju kao prostorni senzori koji mjere okolinu relativno u odnosu na vozilo \cite{suligoj2025percepcija1}. Radi sigurnosti potrebna je redundancija senzora \cite{hrgetic2025uvod}. Glavne modalnosti su:

\begin{itemize}
  \item \textbf{Kamera (2D i 3D)} — bilježi vizualnu informaciju; 2D kamera daje intenzitetsku sliku, dok 3D (stereo) kamera rekonstruira dubinu triangulacijom. Pruža najbogatiju semantičku informaciju, ali je vrlo osjetljiva na osvjetljenje i vremenske uvjete;
  \item \textbf{LiDAR} — laserskim impulsima (mjerenje vremena leta) generira gust 3D oblak točaka visoke geometrijske točnosti; pouzdan noću, ali mu performanse degradiraju raspršenjem u gustoj kiši, magli ili snijegu;
  \item \textbf{Radar} — emitira radiovalove i mjeri odraze; vremensko kašnjenje daje udaljenost, a Dopplerov pomak relativnu brzinu. Robustan je u nepovoljnim uvjetima (kiša, magla, snijeg) i neovisan o osvjetljenju, ali ima nižu prostornu rezoluciju te slabiju osjetljivost na nemetalne objekte;
  \item \textbf{Ultrazvučni senzori} — kratak doseg (do nekoliko metara), jeftini, pogodni za parkiranje i detekciju prepreka pri malim brzinama;
  \item \textbf{GNSS} — globalna satelitska navigacija; uz RTK i IMU postiže centimetarsku točnost lokalizacije;
  \item \textbf{IMU i odometrija kotača} — mjere ubrzanje, kutnu brzinu odnosno rotaciju kotača za procjenu inkrementalnog gibanja vozila.
\end{itemize}

Različite tvrtke biraju različite senzorske strategije \cite{suligoj2025percepcija1}: Waymo i Zoox koriste kombinaciju više kamera, LiDAR-a i radara s naglaskom na redundanciju i točnost, dok Tesla u novijim verzijama koristi pristup temeljen isključivo na kamerama (procjena dubine učenjem iz gibanja), čime smanjuje cijenu hardvera uz veći oslon na podatkovno vođene modele.

\subsection{Fuzija senzora}

Budući da nijedan pojedinačni senzor ne pruža pouzdanu percepciju u svim uvjetima, suvremena vozila oslanjaju se na fuziju senzora koja kombinira komplementarne prednosti kamera, LiDAR-a i radara \cite{suligoj2025percepcija1}. Strategije fuzije obično se dijele prema razini apstrakcije na kojoj se podaci spajaju \cite{suligoj2025percepcija2}:
\begin{enumerate}
  \item \textbf{Rana fuzija} — kombinira sirove podatke prije izlučivanja značajki; omogućuje duboko međumodalno učenje, ali zahtijeva preciznu prostorno-vremensku poravnatost i veliku računsku snagu;
  \item \textbf{Srednja fuzija} — izlučuje značajke specifične za svaki modalitet, a zatim ih spaja u zajedničkom prostoru (često perspektiva iz ptičje perspektive, BEV); uravnotežuje učinkovitost i robusnost;
  \item \textbf{Kasna fuzija} — spaja izlaze pojedinih senzora na razini odluke; čuva modularnost i računski je učinkovita, ali nema dubinske interakcije između modaliteta.
\end{enumerate}
Fuzija također poboljšava procjenu gibanja: kombinacije poput vizualno-inercijske (VIO) ili LiDAR-inercijske odometrije (LIO) smanjuju drift i povećavaju pouzdanost u zahtjevnim uvjetima \cite{suligoj2025percepcija1}.

\subsection{Izazovi percepcije i veza sa sigurnošću}

Vizualna percepcija autonomnih vozila suočava se s nizom izazova: dinamička i složena okruženja, promjenjivi uvjeti osvjetljenja i vremena, zaklanjanja (okluzije), šum i nesigurnost senzora, velike brzine vozila te ograničenja obrade u stvarnom vremenu \cite{suligoj2025percepcija1}. Posebno je ilustrativna ovisnost performansi senzora o okolišnim uvjetima, prikazana u Tablici~\ref{tab:senzori-uvjeti}: dok kamere pružaju najbogatiju informaciju, u magli, jakoj kiši i pri zasljepljujućem suncu njihove performanse značajno opadaju; LiDAR je otporniji, ali pati od raspršenja u oborinama; radar je najrobustniji u nepovoljnim uvjetima, no niske je prostorne rezolucije.

\begin{table}[H]
\centering
\caption{Performanse senzorskih modalnosti u različitim okolišnim uvjetima}
\label{tab:senzori-uvjeti}
\begin{tabularx}{\textwidth}{l X X X X}
\toprule
\textbf{Senzor} & \textbf{Magla} & \textbf{Jaka kiša} & \textbf{Zasljepljujuće sunce} & \textbf{Noć} \\
\midrule
Kamera (mono) & vrlo slabo & slabo & slabo & srednje \\
\addlinespace
Stereo kamera & vrlo slabo & slabo & slabo & srednje \\
\addlinespace
LiDAR & srednje/slabo & srednje & izvrsno & izvrsno \\
\addlinespace
Radar & izvrsno & izvrsno & izvrsno & izvrsno \\
\bottomrule
\end{tabularx}
\end{table}

Upravo ti uvjeti — koji degradiraju percepciju iako senzor i algoritam rade prema specifikaciji — predstavljaju tipične \textit{okidajuće uvjete} u smislu norme ISO~21448 (vidi poglavlje~\ref{sec:sotif}). Stoga su arhitektura percepcije, izbor senzora i njihova fuzija ne samo inženjerska, nego i temeljna sigurnosna pitanja koja povezuju tehničku izvedbu ADS-a s normativnim okvirom obrađenim u nastavku.
